\section{Flujo de trabajo y simulacros de incidentes}

\subsection{Simulacro antes del despliegue}

Antes de cada deployment, debe simularse una ronda de tarea representativa: desde dataset, signature, optimizer hasta los resultados de evidence. El objetivo del simulacro es asegurar que el equipo de DevOps y el prompt obtenido por el optimizer se mantengan alineados.

El proceso del simulacro conviene realizarlo en un entorno de staging: se sube el evidence snapshot al dashboard, para que el equipo de QA verifique de antemano el metric de cada module. Solo cuando el evidence run coincide con el production run (con un error < 5 %) se permite pasar a producción.

\subsection{Proceso de respuesta a incidentes}

El proceso de respuesta a incidentes consta de tres pasos: 1) trace → 2) rollback → 3) annotate. La etapa de trace debe localizar con rapidez qué module produjo el drift; la etapa de rollback debe fijar el last-known-good prompt; la etapa de annotate debe registrar el evento en el governance log.

Al reproducir el trace debe adjuntarse el timeline (Action Timeline / Evidence Matrix) y compararse las marcas de tiempo del metric drift. Esto ayuda al equipo a entender si el drift proviene de la entrada del optimizer o de la salida del tool.

\begin{importantbox}{Checklist del simulacro de incidentes}
\begin{itemize}
    \item Confirmar la condición que dispara la alarm (latency/metric/feedback drift)
    \item Pausar el optimizer y conservar el evidence snapshot
    \item Registrar los recovery steps en el Action Timeline, incluyendo owner y timestamp
\end{itemize}
\end{importantbox}

\subsection{Resumen de la sección}

El flujo de trabajo y los simulacros de incidentes llevan el DSPy pipeline del nivel de research al de producción, y mediante el checklist y el Action Timeline permiten al equipo recuperarse con rapidez cuando ocurre un incidente.
